\documentclass[12pt]{article}
\usepackage[utf8]{inputenc}
\usepackage{amsmath, amssymb, amsfonts}
\usepackage{geometry}
\usepackage{booktabs}
\usepackage{enumitem}
\geometry{margin=1in}
\setlist[enumerate]{leftmargin=*}

\title{Solutions - Practice Problem Set 5\\ \vspace{15pt} \large ECON 101 -- Macroeconomics}
\author{}
\date{Fall 2025}

\begin{document}
\maketitle

\section*{Part A – Multiple Choice Solutions}

\textbf{1. Correct Answer: (C)}  
A lower price level increases consumption, investment, and net exports — causing a movement \emph{down} along the AD curve.



\textbf{2. Correct Answer: (C)}  
A depreciation makes Canadian exports cheaper and imports more expensive, increasing net exports and shifting AD right.



\textbf{3. Correct Answer: (B)}  
A higher price level increases firms’ profits only in the short run, causing a \emph{movement along} SRAS, not a shift.



\textbf{4. Correct Answer: (D)}  
Money multiplier:  
\[
m = \frac{1}{rr} = \frac{1}{0.10} = 10
\]
A new \$10,000 deposit leads to a maximum increase of:  
\[
\Delta M = 10{,}000 \times 9 = 90{,}000
\]



\textbf{5. Correct Answer: (A)}  
An open-market purchase increases reserves, increases the money supply, and lowers the overnight interest rate.


\newpage
\section*{Part B – Long Analytical Problem Solutions}

\textbf{6. Monetary Policy and AD–AS Interaction}

\subsection*{(a) Short-run effects in the AD–AS model}

A drop in consumer confidence reduces autonomous consumption.  
This shifts the AD curve leftward:

- Real GDP falls below potential GDP (\(Y < Y^*\)).  
- The price level falls due to reduced demand.  
- A recessionary gap opens.

---

\subsection*{(b) Change in GDP with MPC = 0.8}

Initial shock:  
\[
\Delta A = -100
\]

Multiplier:  
\[
k = \frac{1}{1 - MPC} = \frac{1}{1 - 0.8} = 5
\]

Total effect on GDP:  
\[
\Delta Y = k \cdot \Delta A = 5 \times (-100) = \boxed{-500}
\]

GDP falls by \$500 billion in total.

---

\subsection*{(c) Bank of Canada policy}

To close the recessionary gap, the Bank of Canada conducts an  
\textbf{open-market purchase of government securities}.  
This:

- Increases bank reserves  
- Expands the money supply  
- Lowers interest rates  
- Stimulates consumption and investment  
- Shifts AD rightward back to potential GDP

---

\subsection*{(d) Required OMO size with money multiplier = 5}

Let \(\Delta R\) = change in reserves from the Bank of Canada's purchase.

Money multiplier:  
\[
m = 5
\]

We need to offset the fall in spending:  
\[
\Delta A = +100
\]

To prevent the multiplied effect of \(-500\), we need:  
\[
\Delta Y = +500
\]

Since \(\Delta M = m \cdot \Delta R\), and an increase in money stimulates AD one-for-one in this simplified setup:

\[
\Delta M = 500
\]

Thus:

\[
\Delta R = \frac{500}{5} = \boxed{100}
\]

The Bank of Canada needs a **\$100 billion open-market purchase** of bonds.

---

\begin{center}
\hrule
\vspace{0.8em}
\textit{End of Instructor Solution Key}
\end{center}

\end{document}
